% Verified artifact-control results; authors must inspect logs and approve. NOT natural clinical duplicate claims.
\subsection{A controlled common-mode confounding challenge}
We evaluated whether high numerical correlation and a shared time
offset necessarily imply common channel-specific digital recording
content. Six distinct public CHB--MIT EDF digital source prefixes
each supplied eight known controlled source transformations and six
negative controls. The additional negative control combines
channel-specific signals from two different published EDF records
with the \emph{same deliberately injected high-amplitude additive
sinusoidal nuisance} at all eight channels. Thus these are known
numerical counterexamples rather than naturally recorded artifacts;
their signal-level labels do not establish physical acquisition or
clinical participant identity. We retain six independent source
units rather than treating the resulting 84 variants as
84 independent observed EEG recordings.

A match-first, verify-second numerical witness required at least
four distinct channels and one shared lag. To reject a common-mode-only
match, we subtracted each time point's mean across the selected
\emph{matched channels} separately from the source and target held-out
verification blocks, then required every channel-specific residual
correlation to exceed the frozen threshold. Selection of the common lag
and assignment used only the original training part and was
\emph{not re-optimised} using the residual verification part.

On these controlled cases, both the raw shared-clock baseline
and the channel-residual witness detected all 48/48 transformed
positives. The uncorrected shared-clock matcher accepted six
of 36 constructed negative cases, all six from the purposely
common-artifact case. The residual witness accepted none of 36.
The public GitHub experiment ran successfully twice, with inspectable
case-level receipts at
\url{https://github.com/lindicaphxag-tech/kaggle/actions/runs/38673797844}
and
\url{https://github.com/lindicaphxag-tech/kaggle/actions/runs/38674323102}.
This establishes the expected behavior of the proposed
\emph{conditional numerical evidence} under one severe constructed
confound, not naturally observed clinical noise immunity.

In particular, subtracting a common channel reference is a classical
EEG preprocessing operation. A fully informed comparator that performs
common average referencing (CAR) before existing shared-clock matching
could remove the same nuisance. Such a comparator, as well as published
multichannel device synchronisation, signal fingerprints and nonlinear
temporal alignment, is required before claiming novel algorithmic
superiority. These controls are not replaced by the positive results
above. At this author-review stage, the residual variant is retained as
a testable research extension, not declared a superior final method.
